\begin{abstract}
Many chat models that run on phones and laptops are 4-bit copies of a model that was evaluated at full precision, and audits of such copies disagree about what quantization does to refusal. A 4-bit build is one of many weight perturbations of the same size, so a change in its behaviour could come from the rounding error specifically or from any perturbation that large. We test this in three 3--4B chat models, with a pre-registered analysis, paired tests, blind human labels and an 8-bit null, and find two 4-bit shifts. With the first tokens of a compliant reply prefilled, Qwen2.5-3B continues into harmful content more often than its original. On held-out jailbreak prompts, Llama-3.2-3B complies \N{C2.cells.llama.4bit-g64.absdiff} points less often, and in an exploratory check it also refuses more benign requests written in jailbreak style. We then rebuild each model with synthetic perturbations matched in size to its rounding error. Shuffled and Gaussian versions shift behaviour in either direction and average close to zero. Scaling the real error, from its mirror image to twice its size, changes behaviour near-linearly over that range. At half the real error, Llama's likelihood is almost unchanged, but it already complies \N{C3.llamafree.e2-scale0.5.absdiff} points less often. In both models the real error moves activations along the refusal direction, but so do builds whose behaviour does not change. Other quantizers at the same bits per weight reproduce the shifts in some cases only, and llama.cpp's Q4\_K\_M shifts neither model. In these models the effect of 4-bit quantization on refusal depends on the specific rounding error and is not predicted by bits per weight or likelihood, so each build should be compared with its original, alongside a null.
\end{abstract}
